\documentclass[../../../../main.tex]{subfiles}
\begin{document}

\begin{flushright}
\footnotesize\textit{【自動文字起こし・要確認】}
\end{flushright}

\noindent
[\textbf{解}]
\begin{enumerate}
    \item[(1)] $f(x) = x^3 + x - 8$ とおく。$f'(x) = 3x^2 + 1 > 0$ から、$f(x)$ は単調増加。
    これと $f(1) = -6$, $f(2) = 2$ より、$f(x) = 0$ は $1 < x < 2$ の区間に唯1つの実根を持つ。

    \item[(2)] 根を $\alpha$ とおくと
    \[
    \alpha^3 + \alpha - 8 = 0 \quad \dots \text{①}
    \]
    $\alpha$ が有理数だと仮定すると、$\alpha = \frac{b}{a}\ (a, b \in \mathbb{N})$ とおける。($(a, b)$ は互いに素) ①に代入して整理し
    \[
    b^3 = 8a^3 - b a^2 = a^2(8a - b)
    \]
    より、$b^3$ が $a^2$ で割り切れることが必要だが、$a, b$ は互いに素だからそのようなことはない。
    よって矛盾し、$\alpha \notin \mathbb{Q}$ である。
\end{enumerate}

\newpage

\end{document}
